% GENERATED FILE. Edit canonical lessons, then run npm run generate:books.
% canonical-source-hash: fnv1a64:102cd9ac93e7185b
% canonical-lessons: FR-C113-oreiller, FR-C113-couverture, FR-C113-serviette, FR-C113-savon, FR-C113-brosse-a-dents

\chapter{Pillow, Blanket, Towel, Soap, Toothbrush}
\label{ch:fr-pillow-blanket-towel-soap-toothbrush}
\hlchaptermodality{113}

\begin{chapteropening}
\textbf{By the end of this chapter:} \emph{I can name a pillow, a blanket, a towel, soap, a toothbrush.}

Complete the last lesson of chapter 113: I can name a pillow, a blanket, a towel, soap, a toothbrush.
\end{chapteropening}

\section[l'oreiller]{l'oreiller — a pillow}

\label{lesson:FR-C113-oreiller}



\begin{quote}
Before the new one: say the French for a shelf, then the French for a curtain.
\end{quote}

\subsection*{You'll want to know: l'oreiller}
\textbf{l'oreiller} — \textquotedblleft{}a pillow\textquotedblright{}.

\begin{cousinweb}
\textbf{oreiller} $\leftarrow$ Old French \emph{orillier} (EtymDB 2.1), built from the word for \textquotedblleft{}ear\textquotedblright{}, $\leftarrow$ Latin \emph{auricula}, \textquotedblleft{}a little ear\textquotedblright{}, with \emph{-ier}: the thing for the ear.
\end{cousinweb}

\begin{grammarlens}[title={What you've built}]
One more word for this chapter.
\end{grammarlens}

\subsection*{Your turn}
\begin{itemize}
\raggedright
  \item \emph{Say it:} \emph{l'oreiller}
  \item \emph{Say it:} \emph{l'oreiller}, once more
  \item \emph{Say it:} say \emph{l'oreiller}
  \item \emph{Recall:} say the French for a shelf, then the French for a curtain, then say \emph{l'oreiller} again
\end{itemize}

\begin{tcolorbox}[breakable,colback=teal!4,colframe=teal!35!black,title={Before you move on}]
What does l'oreiller mean? (\textquotedblleft{}A pillow\textquotedblright{}.) Say it once more.
\end{tcolorbox}
\section[la couverture]{la couverture — a blanket}

\label{lesson:FR-C113-couverture}



\begin{quote}
Before the new one: say the French for a curtain, then the French for a pillow.
\end{quote}

\subsection*{You'll want to know: la couverture}
\textbf{la couverture} — \textquotedblleft{}a blanket\textquotedblright{}.

\begin{cousinweb}
\textbf{couverture} $\leftarrow$ Late Latin \emph{coopertura}, \textquotedblleft{}a covering\textquotedblright{}, from \emph{coopertus}, \textquotedblleft{}covered\textquotedblright{}, the past participle of \emph{cooperire} — inherited. (EtymDB 2.1)
\end{cousinweb}

\begin{grammarlens}[title={What you've built}]
One more word for this chapter.
\end{grammarlens}

\subsection*{Your turn}
\begin{itemize}
\raggedright
  \item \emph{Say it:} \emph{la couverture}
  \item \emph{Say it:} \emph{la couverture}, once more
  \item \emph{Say it:} say \emph{la couverture}
  \item \emph{Recall:} say the French for a curtain, then the French for a pillow, then say \emph{la couverture} again
\end{itemize}

\begin{tcolorbox}[breakable,colback=teal!4,colframe=teal!35!black,title={Before you move on}]
What does la couverture mean? (\textquotedblleft{}A blanket\textquotedblright{}.) Say it once more.
\end{tcolorbox}
\section[la serviette]{la serviette — a towel; a napkin}

\label{lesson:FR-C113-serviette}



\begin{quote}
Before the new one: say the French for a pillow, then the French for a blanket.
\end{quote}

\subsection*{You'll want to know: la serviette}
\textbf{la serviette} — \textquotedblleft{}a towel; a napkin\textquotedblright{}.

\begin{cousinweb}
\textbf{serviette} — built from the verb for \textquotedblleft{}to serve\textquotedblright{}, with \emph{-ette}: a little thing for serving. The verb $\leftarrow$ Latin \emph{servire}, \textquotedblleft{}to be a servant\textquotedblright{}. (EtymDB 2.1)
\end{cousinweb}

\begin{grammarlens}[title={What you've built}]
One more word for this chapter.
\end{grammarlens}

\subsection*{Your turn}
\begin{itemize}
\raggedright
  \item \emph{Say it:} \emph{la serviette}
  \item \emph{Say it:} \emph{la serviette}, once more
  \item \emph{Say it:} say \emph{la serviette}
  \item \emph{Recall:} say the French for a pillow, then the French for a blanket, then say \emph{la serviette} again
\end{itemize}

\begin{tcolorbox}[breakable,colback=teal!4,colframe=teal!35!black,title={Before you move on}]
What does la serviette mean? (\textquotedblleft{}A towel; a napkin\textquotedblright{}.) Say it once more.
\end{tcolorbox}
\section[le savon]{le savon — soap}

\label{lesson:FR-C113-savon}



\begin{quote}
Before the new one: say the French for a blanket, then the French for a towel; a napkin.
\end{quote}

\subsection*{You'll want to know: le savon}
\textbf{le savon} — \textquotedblleft{}soap\textquotedblright{}.

\begin{cousinweb}
\textbf{savon} $\leftarrow$ Old French \emph{savon} (EtymDB 2.1) $\leftarrow$ Latin \emph{sapo}, \emph{saponis}, \textquotedblleft{}soap\textquotedblright{}, a Germanic word in Latin (Lewis \& Short) — inherited.
\end{cousinweb}

\begin{grammarlens}[title={What you've built}]
One more word for this chapter.
\end{grammarlens}

\subsection*{Your turn}
\begin{itemize}
\raggedright
  \item \emph{Say it:} \emph{le savon}
  \item \emph{Say it:} \emph{le savon}, once more
  \item \emph{Say it:} say \emph{le savon}
  \item \emph{Recall:} say the French for a blanket, then the French for a towel; a napkin, then say \emph{le savon} again
\end{itemize}

\begin{tcolorbox}[breakable,colback=teal!4,colframe=teal!35!black,title={Before you move on}]
What does le savon mean? (\textquotedblleft{}Soap\textquotedblright{}.) Say it once more.
\end{tcolorbox}
\section[la brosse à dents]{la brosse à dents — a toothbrush}

\label{lesson:FR-C113-brosse-a-dents}



\begin{quote}
Before the new one: say the French for a towel; a napkin, then the French for soap.
\end{quote}

\subsection*{You'll want to know: la brosse à dents}
\textbf{la brosse à dents} — \textquotedblleft{}a toothbrush\textquotedblright{}.

\begin{cousinweb}
\textbf{brosse à dents} — \textquotedblleft{}brush for teeth\textquotedblright{}: \emph{brosse}, \textquotedblleft{}brush\textquotedblright{}, first \textquotedblleft{}brushwood\textquotedblright{} in Old French (EtymDB 2.1), of uncertain origin; \emph{à}, \textquotedblleft{}for\textquotedblright{}; and \emph{dents}, \textquotedblleft{}teeth\textquotedblright{}, $\leftarrow$ Latin \emph{dens}, \emph{dentis}.
\end{cousinweb}

\begin{grammarlens}[title={What you've built}]
One more word for this chapter.
\end{grammarlens}

\subsection*{Your turn}
\begin{itemize}
\raggedright
  \item \emph{Say it:} \emph{la brosse à dents}
  \item \emph{Say it:} \emph{la brosse à dents}, once more
  \item \emph{Say it:} say \emph{la brosse à dents}
  \item \emph{Recall:} say the French for a towel; a napkin, then the French for soap, then say \emph{la brosse à dents} again
\end{itemize}

\begin{tcolorbox}[breakable,colback=teal!4,colframe=teal!35!black,title={Before you move on}]
What does la brosse à dents mean? (\textquotedblleft{}A toothbrush\textquotedblright{}.) Say it once more.
\end{tcolorbox}
